% ---------------------------------------------------------------------------
\section{Receiver pressure: nonlinear contact and collapse-wave estimates}\label{sec:receiver}
Handoff priority 6. The pressure observables O5--O7 are computed here from a stated jet speed; they are
never identified with the source pressure O4.

\subsection{Step P1 --- Weak-shock Hugoniot of water closed with the Tait equation}
\emph{Symbols} as in Limits E11: upstream state $(\rho_0,p_0)$ at rest, shocked state $(\rho_1,p_1)$, shock speed
$D$ relative to the upstream liquid, particle-velocity jump $u_p$. Tait isentrope for water (Lechner et al.\ 2020,
Eq.~8, verified): $n_{\mathrm T}=7.15$, $B_{\mathrm T}=305$\,MPa.
Here $\rho_0,\rho_1>0$ are densities (kg/m$^3$), $p_0,p_1$ absolute pressures (Pa),
$D>u_p\ge0$ shock and particle-jump speeds (m/s), $n_{\mathrm T}$ dimensionless exponent,
and $B_{\mathrm T}$ pressure offset (Pa); subscript 0 denotes upstream and 1 compressed liquid:
\begin{equation}
\rho_0D=\rho_1(D-u_p),\qquad p_1-p_0=\rho_0Du_p,\qquad p_1=(p_0+B_{\mathrm T})\Bigl(\frac{\rho_1}{\rho_0}\Bigr)^{n_{\mathrm T}}-B_{\mathrm T}.\label{eq:hug}
\end{equation}
The first two relations are exact jump conditions (mass and momentum, Limits E11). The third replaces the
Hugoniot by the isentrope; the two curves osculate to second order in the specific-volume jump, the entropy
jump being third order (Landau \& Lifshitz, \emph{Fluid Mechanics}, \S86), so the closure is accurate for weak
to moderate shocks and is flagged above 1\,GPa. Eliminating $\rho_1$ gives one equation for $D$:
$(p_0+B_{\mathrm T})[D/(D-u_p)]^{n_{\mathrm T}}-B_{\mathrm T}-p_0-\rho_0Du_p=0$, solved by bracketing.
\emph{Limit check:} as $u_p\to0$, $D\to c_{\mathrm T}=\sqrt{n_{\mathrm T}(p_0+B_{\mathrm T})/\rho_0}=\cTait$\,m/s and
$p_1-p_0\to\rho_0c_{\mathrm T}u_p$, recovering Limits E09 (rigid limit).

\paragraph{Rigid receiver.} Arrest of a normally arriving plug of speed $U$ is the piston problem with $u_p=U$.
\paragraph{Linear-impedance receiver.} Here $U$ is incident jet speed (m/s), $v$ interface speed
(m/s), $Z_t\ge0$ opposing impedance (Pa\,s/m), and $p_{\mathrm{Hug}}$ the jet-side
Tait-closed absolute pressure (Pa). The mixed closure solves
$p_{\mathrm{Hug}}(U-v)-p_0=Z_tv$. It treats the jet side nonlinearly and the opposing medium
linearly. Setting $Z_t=Z_{\mathrm w}=\rho_0c_{\mathrm T}$ estimates an equal-impedance contact scale;
it is not the symmetric nonlinear two-water-shock solution. At low Mach number both give
$p_{\mathrm c}\approx\tfrac12Z_{\mathrm w}U$
at low Mach number --- the factor $\tfrac12$ of Supponen et al.\ (2017, Eq.~4, verified:
$p_h=\tfrac12\rho cU_{\mathrm{jet}}=0.45(\rho c^2\Delta p)^{1/2}\zeta^{-1}$, stated valid for $\zeta>10^{-3}$ and
failing near $\zeta\lesssim0.9(\Delta p/\rho)^{1/2}/c$).

\input{generated/table_receiver.tex}

\textbf{Results.} Table~\ref{tab:receiver}. The nonlinear first-contact pressure exceeds the linear water-hammer
value by $\ratioUhundred\times$ at 100\,m/s, $\ratioUthreehundred\times$ at 300\,m/s and $\ratioUeightfifty\times$ at
850\,m/s; Limits E10 is therefore \emph{not} an upper bound outside the small-Mach regime. At the prior sampled
outlet speed 25.28\,m/s the Tait value is $\pcUfivetwenty$\,MPa (linear: 37.9\,MPa).

\subsection{Step P2 --- Duration and area: why a contact pressure is not a measured peak}
The one-dimensional hammer pressure persists only until the lateral release wave from the jet edge reaches
the axis, $t_{\mathrm r}\approx a_{\mathrm j}/c_1$ with $a_{\mathrm j}$ the jet radius and $c_1$ the shocked sound
speed. For a jet of radius $1\,\mu$m, $t_{\mathrm r}\approx0.6$\,ns: a sensor with 30\,MHz bandwidth (15\,ns rise
time, Supponen et al.\ 2017) averages the contact pressure over an interval about 25 times longer and an area
much larger than the jet, so the measured peak (O7) is far below O6. Any comparison with measurements must
apply the sensor's area and impulse response (Limits E12--E13, E32).

\subsection{Step P3 --- Benchmark B1: what loads the receiver}
For $\gamma\ge1.25$ the jet first strikes the opposite cavity wall, not the receiver. The internal impact is a
liquid--liquid contact (equal impedance) producing a local hammer pressure scale
$p_h\approx\tfrac12Z_{\mathrm w}U_{\mathrm j}$. The mixed contact calculation above supplies the
nonlinear/linear estimate tabulated in Table~\ref{tab:envelope}; it is not a resolved internal
shock simulation. That contact emits a torus-like wave from an assumed jet-tip area.
The receiver, a distance $d\approx z_{\mathrm{bot}}$ below the impact point, is reached
by this wave. Under linear spherical spreading from a source of radius $a_{\mathrm j}$, and rigid-wall reflection
of an incident wave (factor 2, Limits E14), the receiver pressure is estimated as follows.
Here $p_{\mathrm{rec}}$ is a receiver-wave peak estimate (Pa), $p_h$ the internal
liquid--liquid contact increment (Pa), $a_{\mathrm j}>0$ source/jet radius (m), and $d>0$
the propagation separation (m). The minimum caps the spreading factor at unity.
This is an estimated extension of wave observable O5 to a receiver plane; it is not direct
liquid contact O6 or sensor-averaged O7:
\begin{equation}
p_{\mathrm{rec}}\approx2\,p_h\,\min\Bigl(1,\frac{a_{\mathrm j}}{d}\Bigr).\label{eq:prec}
\end{equation}
This is an order-of-magnitude estimate (the wave is not spherical, the near field is nonlinear and the cavity
remnant scatters it); it is reported as such in Table~\ref{tab:envelope} and is not a resolved receiver pressure.
Direct jet contact with the receiver requires $\gamma\lesssim1$, outside $\mathcal D$, where the cavity touches the wall
and potential flow fails; that case remains unresolved (Section~\ref{sec:validity}).

\subsection{Step P4 --- Conditional bound for a declared pressure observer}
The accepted progressive-wave relations of Limits E12--E14 give a pressure bound only after
fixing the observer area and time scale. Assume a locally planar, normally propagating
\emph{linear} wave with carrier impedance $Z_{\mathrm w}=\rho_{\mathrm w}c_{\mathrm w}>0$
(Pa\,s/m), incident pressure increment $p_{\mathrm{inc}}(\mathbf x,t)$ (Pa), position $\mathbf x$ (m),
time $t$ (s), observer region
$\Omega_{\mathrm o}$ of area $A_{\mathrm o}>0$ (m$^2$), and interval $\tau_{\mathrm o}>0$ (s).
Define $E_{A\tau}$ as the incident acoustic energy crossing that region in that interval (J),
$\bar p_{\mathrm{inc}}$ its area/time mean (Pa), and $p_{\mathrm{rms}}$ its root mean square (Pa).
Area/time elements $\dd A,\dd t$ have units m$^2$, s, and $\Phia\ge0$ is activated exergy (J).
The wave-energy flux is $p_{\mathrm{inc}}^2/Z_{\mathrm w}$; with the no-external-work budget
of Step T1, the captured energy cannot exceed the total available energy:
\begin{align}
E_{A\tau}&=\frac{1}{Z_{\mathrm w}}\int_0^{\tau_{\mathrm o}}\int_{\Omega_{\mathrm o}}
p_{\mathrm{inc}}^2\,\dd A\,\dd t\le\Phia,\notag\\
p_{\mathrm{rms}}^2&=\frac{1}{A_{\mathrm o}\tau_{\mathrm o}}
\int_0^{\tau_{\mathrm o}}\int_{\Omega_{\mathrm o}}p_{\mathrm{inc}}^2\,\dd A\,\dd t
\le\frac{Z_{\mathrm w}\Phia}{A_{\mathrm o}\tau_{\mathrm o}}.\label{eq:pressure_rms}
\end{align}
Applying Cauchy--Schwarz to $p_{\mathrm{inc}}$ and the constant function 1 on the same
area/time region gives $(\int\!\int p_{\mathrm{inc}})^2\le
(A_{\mathrm o}\tau_{\mathrm o})\int\!\int p_{\mathrm{inc}}^2$, hence
$|\bar p_{\mathrm{inc}}|\le p_{\mathrm{rms}}$.
For a linear opposing medium with constant impedance $Z_t>0$ (Pa\,s/m), define the pressure
transmission factor $\mathcal T_p=2Z_t/(Z_t+Z_{\mathrm w})\in(0,2)$ (dimensionless), and
$\bar p_{\mathrm{rec}}$ as the corresponding area/time mean at the receiver (Pa):
\begin{equation}
|\bar p_{\mathrm{rec}}|\le\mathcal T_p
\sqrt{\frac{Z_{\mathrm w}\Phia}{A_{\mathrm o}\tau_{\mathrm o}}}
\le2\sqrt{\frac{Z_{\mathrm w}\Phia}{A_{\mathrm o}\tau_{\mathrm o}}}.
\label{eq:pressure_observer}
\end{equation}
The factor 2 is the rigid-reflection limit, not an additional factor for an already calculated
collision pressure. Units check: $Z_{\mathrm w}\Phia/(A_{\mathrm o}\tau_{\mathrm o})$ is Pa$^2$.
The incident bound is saturated algebraically by a uniform fixed-sign pulse with all the budget
assigned to that wave; this does not prove that PFC collapse can generate that pulse.
As $A_{\mathrm o}\tau_{\mathrm o}\to0$ the bound grows without supplying a finite point-pressure
cap. These are conditional linear-wave/RMS or averaged-pressure bounds, not bounds on an
instantaneous nonlinear contact shock. An actual sensor needs its spatial and temporal response.
